\documentclass[11pt]{article}
\usepackage[margin=1in]{geometry}
\usepackage{amsmath,amssymb,amsthm,mathtools}
\usepackage{graphicx}
\usepackage{xcolor}
\usepackage[hidelinks]{hyperref}

\newtheorem{theorem}{Theorem}
\newtheorem{lemma}{Lemma}
\newtheorem{corollary}{Corollary}
\newcommand{\BH}{B(\mathcal H)}
\newcommand{\Sone}{S_1(\mathcal H)}

\title{Normal maps on $L^\infty$ induced by operator-valued measures\\
\large A full answer to Problem 2 of arXiv:1411.4559}
\author{Automated mathematics run \texttt{fa\_banach\_001}, lane 13}
\date{22 July 2026}

\begin{document}
\maketitle

\noindent
\fcolorbox{black}{gray!10}{\parbox{0.94\textwidth}{
\textbf{Status.} Claimed full solution to the literal Problem 2, likely valid
subject to human review.  An inducing ultraweakly continuous linear map exists
if and only if the operator-valued measure vanishes on all $\mu$-null sets.
Without this condition the answer is negative, even for a probability
operator-valued measure on a one-dimensional Hilbert space.
}}

\section{Source problem}

The source is D. Han, D. R. Larson, B. Liu, and R. Liu,
\emph{Dilations of frames, operator valued measures and bounded linear maps},
arXiv:1411.4559v2, Problem 2 on PDF page 19 \cite{HLLL}.
The paper defines an operator-valued measure
$E:\Sigma\to B(\mathcal H)$ to be countably additive in the weak operator
topology.  Its question is:

\begin{quote}
Given a probability space $(\Omega,\Sigma,\mu)$ and an operator-valued measure
$E:(\Omega,\Sigma)\to B(\mathcal H)$, does there exist an ultraweakly
continuous map $\phi:L^\infty(\mu)\to B(\mathcal H)$ such that
$E(B)=\phi(\chi_B)$ for every $B\in\Sigma$?  If not, determine necessary and
sufficient conditions.
\end{quote}

\begin{figure}[ht]
\centering
\includegraphics[width=\textwidth]{figures/open_problem_crop.png}
\caption{Problem 2 in the source paper, PDF page 19.}
\end{figure}

\section{Definitions}

Write $E\ll\mu$ when
\[
 \mu(B)=0\quad\Longrightarrow\quad E(B)=0
 \qquad(B\in\Sigma).
\]
We use the canonical preduals $L^1(\mu)$ of $L^\infty(\mu)$ and
$\Sone$ of $\BH$.  Thus ``ultraweakly continuous'' means weak-star to
weak-star continuous, or equivalently that the map is the adjoint of a bounded
map $\Sone\to L^1(\mu)$.

For a bounded simple function $s=\sum_j a_j\chi_{B_j}$ on a finite measurable
partition, put
\[
 \int s\,dE:=\sum_j a_j E(B_j).
\]
The definition is independent of the chosen representation.

\section{Proof intuition}

The obstruction is forced by the quotient defining $L^\infty(\mu)$: the
indicator of a null set is the zero element, so any inducing map must kill the
operator assigned to that set.  The striking point is that there is no further
obstruction.  Weak-operator countable additivity gives a uniform bound for all
simple integrals.  Absolute continuity then lets scalarizations by trace-class
operators acquire $L^1$ Radon--Nikodym densities.  Those densities form a
bounded preadjoint, and its adjoint is the required normal map.

\section{The characterization}

\begin{lemma}[Uniform semivariation]\label{lem:semi}
Let $E:\Sigma\to\BH$ be weak-operator countably additive.  Then
\[
 C_E:=\sup\left\{\left\|\int s\,dE\right\|:
 s\text{ simple},\ |s|\leq 1\right\}<\infty.
\]
Consequently, for every $\rho\in\Sone$, the scalar set function
\[
 \nu_\rho(B):=\operatorname{Tr}(E(B)\rho)
\]
is a countably additive complex measure and
\[
 |\nu_\rho|(\Omega)\leq C_E\|\rho\|_1.
\]
\end{lemma}

\begin{proof}
For $x\in\mathcal H$, the set function $E_x(B)=E(B)x$ is weakly countably
additive.  By the Orlicz--Pettis theorem it is a norm-countably-additive
$\mathcal H$-valued measure.  Every countably additive Banach-space-valued
measure has finite semivariation \cite[Chapter I]{DU}; hence, for each fixed
$x$,
\[
 \sup_{s\text{ simple},\ |s|\leq1}
 \left\|\left(\int s\,dE\right)x\right\|<\infty.
\]
The uniform boundedness principle applied to the family of simple-integration
operators gives $C_E<\infty$.

Weak-operator countable additivity, together with uniform boundedness of the
range of $E$, implies countable additivity after pairing with any trace-class
operator: approximate that trace-class operator in trace norm by finite-rank
operators.  Thus $\nu_\rho$ is a complex measure.  The standard partition
formula for total variation and trace duality give
\begin{align*}
 |\nu_\rho|(\Omega)
 &=\sup_{\substack{(B_j)\text{ finite partition}\\ |a_j|\leq1}}
 \left|\sum_j a_j\operatorname{Tr}(E(B_j)\rho)\right|\\
 &=\sup_{s\text{ simple},\ |s|\leq1}
 \left|\operatorname{Tr}\left(\left(\int s\,dE\right)\rho\right)\right|\\
 &\leq C_E\|\rho\|_1.
\end{align*}
\end{proof}

\begin{theorem}\label{thm:main}
Let $(\Omega,\Sigma,\mu)$ be a probability space, let $\mathcal H$ be any
Hilbert space, and let $E:\Sigma\to\BH$ be weak-operator countably additive.
The following are equivalent.
\begin{enumerate}
\item There exists an ultraweakly continuous bounded linear map
$\phi:L^\infty(\mu)\to\BH$ satisfying
$\phi(\chi_B)=E(B)$ for every $B\in\Sigma$.
\item $E\ll\mu$.
\end{enumerate}
When these conditions hold, $\phi$ is unique and is given by weak operator
integration,
\[
 \langle\phi(f)x,y\rangle=\int_\Omega f\,d\langle E(\cdot)x,y\rangle.
\]
If $E(\Omega)=I_{\mathcal H}$, then $\phi$ is unital.
\end{theorem}

\begin{proof}
Suppose first that such a $\phi$ exists.  If $\mu(B)=0$, then
$\chi_B=0$ as an element of $L^\infty(\mu)$, and therefore
$E(B)=\phi(\chi_B)=0$.  Hence $E\ll\mu$.

Conversely assume $E\ll\mu$.  For $\rho\in\Sone$, Lemma~\ref{lem:semi}
gives a finite complex measure $\nu_\rho$ with
$|\nu_\rho|(\Omega)\leq C_E\|\rho\|_1$.  Absolute continuity of $E$ implies
$\nu_\rho\ll\mu$.  By the scalar Radon--Nikodym theorem there is a unique
$g_\rho\in L^1(\mu)$ such that
\[
 \nu_\rho(B)=\int_B g_\rho\,d\mu,
 \qquad \|g_\rho\|_1=|\nu_\rho|(\Omega)leq C_E\|\rho\|_1.
\]
Uniqueness of Radon--Nikodym derivatives shows that
\[
 \Psi:\Sone\longrightarrow L^1(\mu),\qquad \Psi(\rho)=g_\rho,
\]
is linear, and the displayed estimate shows that it is bounded.  Define
\[
 \phi:=\Psi^*:L^\infty(\mu)\longrightarrow\BH
\]
using the canonical dualities.  This map is bounded and ultraweakly
continuous by construction.  For $B\in\Sigma$ and $\rho\in\Sone$,
\begin{align*}
 \operatorname{Tr}(\phi(\chi_B)\rho)
 &=\int_B g_\rho\,d\mu\\
 &=\nu_\rho(B)
 =\operatorname{Tr}(E(B)\rho).
\end{align*}
Since trace-class operators separate points of $\BH$, we obtain
$\phi(\chi_B)=E(B)$.

The weak integration formula follows first for simple functions and then for
all of $L^\infty$ by uniform approximation by simple functions.  The same
approximation proves uniqueness.  Finally,
$\phi(1)=\phi(\chi_\Omega)=E(\Omega)$, proving the unital assertion.
\end{proof}

\section{Sharp negative example}

\begin{corollary}
The unrestricted existence question in Problem 2 has a negative answer, even
when $\mathcal H=\mathbb C$ and $E$ is a probability operator-valued measure.
\end{corollary}

\begin{proof}
Take $\Omega=[0,1]$ with Lebesgue measure and let
\[
 E(B)=\mathbf 1_{\{0\in B\}}I_{\mathcal H}.
\]
This is a weak-operator countably additive probability measure, but
$E(\{0\})=I_{\mathcal H}$ while $\chi_{\{0\}}=0$ in $L^\infty(\mu)$.
No inducing map, normal or otherwise, can exist.
\end{proof}

\section{Verification and limitations}

The proof was checked against the source definition of operator-valued
measure and uses no separability assumption.  The main verifier focus should
be Lemma~\ref{lem:semi}: the standard finite-semivariation theorem for
countably additive vector measures, followed by uniform boundedness, supplies
exactly the uniform constant required for the trace-class preadjoint.

A bounded exact-phrase and close-variant search was run on 22 July 2026 using
the source title, ``Problem 2'', ``operator valued measure'',
``ultraweakly continuous map'', and $L^\infty(\mu)$.  It found the source and
related later work on normal \emph{dilations}, but no paper explicitly
answering this induction problem.  This packet therefore claims a full proof
for the literal source question, not priority or publication-level novelty.
The characterization is a short consequence of classical vector-measure
theory and may be folklore.

\begin{thebibliography}{9}
\bibitem{HLLL}
D. Han, D. R. Larson, B. Liu, and R. Liu,
\emph{Dilations of frames, operator valued measures and bounded linear maps},
arXiv:1411.4559v2 (2014), Problem 2, p.~19.

\bibitem{DU}
J. Diestel and J. J. Uhl, Jr.,
\emph{Vector Measures}, Mathematical Surveys, vol.~15,
American Mathematical Society, 1977.
\end{thebibliography}

\end{document}
